\chapter{Formulation}

\section{Places}

\begin{itemize}

\item Let $\mathfrak{X}$ denote a fixed place, which will be understood as a
fundamental and elementary answer to the question of ``where?'' without 
recourse to concepts of location or position, which require at least 
two fixed places to be defined.

\item Let $\mathbb{S}$ denote space, defined as the infinite set of fixed places
$\mathfrak{X}$.

\end{itemize}

\section{Space Coordinates}
\begin{itemize}
\item From the infinite set of fixed places $\mathbb{S}$, we select
one fixed place to designate as the {\bf origin} $O$.

\item At the origin $O$, we construct three dextral, orthonormal {\bf unit
vectors}, $\veeonehat, \veetwohat, \veethreehat$, called the 
{\bf spatial basis vectors}.

\item The term {\bf dextral} designates a {\em right-handed rule}, 
such that 
$\veeonehat \times \veetwohat = \veethreehat$, 
$\veetwohat \times \veethreehat = \veeonehat$, and
$\veethreehat \times \veeonehat = \veetwohat$, or more simply, that
$\veeihat \times \veejhat = \veekhat$.

\item The origin $O$ and the spatial basis vectors $\veeonehat, \veetwohat, 
\veethreehat$ together the {\bf spatial coordinate system} $S$, which is
a one-to-one correspondence between fixed places $\mathfrak{X}$
in space $\mathbb{S}$ and their {\bf spatial coordinates} $\vx$.

\item The spatial coordinates are an ordered set of three real numbers
$\{x_1, x_2, x_3\}\transpose$ measured in a particular coordinate system
from the coordinate system's origin to the fixed place $\mathfrak{X}$
in space $\mathbb{S}$.
\item A place $\mathfrak{X}$ in space $\mathbb{S}$ 
is located the spatial coordinate system $S$ by the 
particular spatial
coordinate $\vx$, written explicitly as
\be
\vx = x_1 \veeonehat + x_2 \veetwohat + x_3 \veethreehat = x_i \veeihat.
\ee
\item Since, at the present moment, 
we have only introduced one coordinate system (and thus only
one set of basis vectors), we need not designate the basis vectors explicitly, 
and may simply write 
\be
\vx = \{x_1, x_2, x_3\}\transpose
\label{eq:no_basis}
\ee 
without ambiguity.

\item This abbreviated form (\ref{eq:no_basis}), 
however, will cause ambiguity
later, when two or more coordinate systems are present to describe
places in space.  In this latter case, 
the explicit basis vectors or the explicit coordinate system must be 
stated to avoid confusion.  

\end{itemize}
